% Laget automatisk av elevark/lag_elevark.py fra sim/myoner.html. Endringer her blir overskrevet.
\documentclass[11pt]{article}
\usepackage{elevark}
\begin{document}
\elevarkhode{Fysikk 2 · Spesiell relativitetsteori}{Myoner fra atmosfæren}{https://fysikklab.no/sim/myoner.html}{Simuleringen}

Når kosmisk stråling treffer atmosfæren, dannes det myoner høyt oppe. Et myon i ro lever i gjennomsnitt bare 2,2~µs før det henfaller. Likevel måler vi mange myoner ved bakken. Her følger vi 1000 myoner på vei ned, med og uten relativitetsteori.

\section*{Oppgaver}

Antallet som er igjen etter tiden $t$, er $N = N_{0} \cdot e^{-t/τ}$, der $τ = 2{,}2\,\mathrm{\symup{\mu}s}$ er den midlere levetiden. Det tilsvarer en halveringstid på 1,5~µs.

\begin{oppgaveliste}
\oppgave \fase{gjett}{Gjett}1000 myoner lages 10~km over bakken og går rett ned med farten $0{,}995c$. Hvor mange når bakken?
\svar
\oppgave \fase{test}{Regn}Regn ut hvor lang tid turen tar sett fra jorda. Hvor mange myoner skulle nådd bakken hvis myonets klokke gikk like fort som våre?
\svar
\oppgave \fase{test}{Test}Kjør simuleringen. Hvordan går det i de to kolonnene?
\svar
\oppgave \fase{forklar}{Forklar}Bruk tidsforlengelse til å forklare hvorfor så mange når bakken. Hvor lang tid har gått på myonets egen klokke når det kommer ned? Regn ut hvor mange som skal nå bakken.
\svar
\oppgave \fase{forklar}{Forklar}Bytt referansesystem til myonene. Nå står myonene i ro, og klokkene deres går normalt. Hvorfor når like mange bakken likevel?
\svar
\oppgave \fase{gjett}{Gjett}Trykk «Nye myoner» og kjør flere ganger. Hvorfor blir ikke antallet ved bakken det samme hver gang?
\svar
\oppgave \fase{utfordring}{Utfordring}I 1963 målte Frisch og Smith myoner både på toppen av et fjell og ved havnivå, ca. 1900~m lenger nede. Sett høyden til 2~km og $v = 0{,}995c$. Hvor stor andel overlever turen med og uten relativitetsteori?
\svar
\end{oppgaveliste}

\end{document}
